\documentclass[11pt]{article}

\usepackage[margin=1in]{geometry}
\usepackage{amsmath,amssymb}
\usepackage{array}
\usepackage{booktabs}
\usepackage{enumitem}
\usepackage[hidelinks]{hyperref}

\setlist[itemize]{leftmargin=1.4em}
\setlist[enumerate]{leftmargin=1.6em}

\title{SYSC 5108 Exam Study\\Assignment 4, Question 8}
\author{Jesse Levine}
\date{}

\begin{document}
\maketitle

\section*{Question}

Let \(X=(X_1,X_2)\) be a two-dimensional random variable with binary entries
\[
    X_1,X_2 \in \{0,1\}.
\]
The joint probabilities are
\[
    P(X_1,X_2)
    =
    \begin{bmatrix}
        P(0,0)\\
        P(0,1)\\
        P(1,0)\\
        P(1,1)
    \end{bmatrix}
    =
    \begin{bmatrix}
        \frac14\\[4pt]
        \frac12\\[4pt]
        \frac18\\[4pt]
        \frac18
    \end{bmatrix}.
\]

\begin{enumerate}[label=(\alph*)]
    \item Find the distributions \(P(X_1\mid X_2)\) and \(P(X_2\mid X_1)\).
    \item Run Gibbs sampling for \(3\) iterations starting from uniform
    distributions for both \(X_1\) and \(X_2\) independently, using the random
    numbers
    \[
        0.94,\ 0.27,\ 0.30,\ 0.18,\ 0.75,\ 0.35,\ 0.61,\ 0.42,\ldots
    \]
\end{enumerate}

\section*{Course and Textbook Cross-References}

\begin{itemize}
    \item \textbf{Chapter 7 slides, Gibbs Sampling:} repeatedly sample one
    variable from its conditional distribution given the other variables.
    \item \textbf{Chapter 7 slides, Sampling for Structured Models:} Gibbs
    sampling uses conditional distributions \(p(x_j\mid x_{-j})\).
    \item \textbf{Goodfellow, Bengio, and Courville, \emph{Deep Learning},
    Chapter 17:} Markov chain Monte Carlo updates the state by repeated sampling
    from conditional distributions.
\end{itemize}

\section*{Step 1: Write the Joint Table Clearly}

The lexicographic ordering means
\[
    P(0,0)=\frac14,\qquad
    P(0,1)=\frac12,\qquad
    P(1,0)=\frac18,\qquad
    P(1,1)=\frac18.
\]

So the joint table is
\[
\begin{array}{c|cc}
    & X_2=0 & X_2=1\\
    \hline
    X_1=0 & \frac14 & \frac12\\[4pt]
    X_1=1 & \frac18 & \frac18
\end{array}
\]

\section*{Part (a): Find \(P(X_1\mid X_2)\) and \(P(X_2\mid X_1)\)}

\subsection*{Step 2: Compute the Marginals}

For \(X_2\),
\[
    P(X_2=0)=P(0,0)+P(1,0)=\frac14+\frac18=\frac38,
\]
\[
    P(X_2=1)=P(0,1)+P(1,1)=\frac12+\frac18=\frac58.
\]

For \(X_1\),
\[
    P(X_1=0)=P(0,0)+P(0,1)=\frac14+\frac12=\frac34,
\]
\[
    P(X_1=1)=P(1,0)+P(1,1)=\frac18+\frac18=\frac14.
\]

\subsection*{Step 3: Compute \(P(X_1\mid X_2)\)}

When \(X_2=0\),
\[
    P(X_1=0\mid X_2=0)=\frac{P(0,0)}{P(X_2=0)}=\frac{\frac14}{\frac38}=\frac23,
\]
\[
    P(X_1=1\mid X_2=0)=\frac{P(1,0)}{P(X_2=0)}=\frac{\frac18}{\frac38}=\frac13.
\]
Hence
\[
    P(X_1\mid X_2=0)
    =
    \begin{bmatrix}
        \frac23\\[4pt]
        \frac13
    \end{bmatrix}.
\]

When \(X_2=1\),
\[
    P(X_1=0\mid X_2=1)=\frac{P(0,1)}{P(X_2=1)}=\frac{\frac12}{\frac58}=\frac45,
\]
\[
    P(X_1=1\mid X_2=1)=\frac{P(1,1)}{P(X_2=1)}=\frac{\frac18}{\frac58}=\frac15.
\]
Hence
\[
    P(X_1\mid X_2=1)
    =
    \begin{bmatrix}
        \frac45\\[4pt]
        \frac15
    \end{bmatrix}.
\]

\subsection*{Step 4: Compute \(P(X_2\mid X_1)\)}

When \(X_1=0\),
\[
    P(X_2=0\mid X_1=0)=\frac{P(0,0)}{P(X_1=0)}=\frac{\frac14}{\frac34}=\frac13,
\]
\[
    P(X_2=1\mid X_1=0)=\frac{P(0,1)}{P(X_1=0)}=\frac{\frac12}{\frac34}=\frac23.
\]
Hence
\[
    P(X_2\mid X_1=0)
    =
    \begin{bmatrix}
        \frac13\\[4pt]
        \frac23
    \end{bmatrix}.
\]

When \(X_1=1\),
\[
    P(X_2=0\mid X_1=1)=\frac{P(1,0)}{P(X_1=1)}=\frac{\frac18}{\frac14}=\frac12,
\]
\[
    P(X_2=1\mid X_1=1)=\frac{P(1,1)}{P(X_1=1)}=\frac{\frac18}{\frac14}=\frac12.
\]
Hence
\[
    P(X_2\mid X_1=1)
    =
    \begin{bmatrix}
        \frac12\\[4pt]
        \frac12
    \end{bmatrix}.
\]

\subsection*{Answer for Part (a)}

\[
    \boxed{
    P(X_1\mid X_2=0)=
    \begin{bmatrix}
        \frac23\\[4pt]
        \frac13
    \end{bmatrix},
    \qquad
    P(X_1\mid X_2=1)=
    \begin{bmatrix}
        \frac45\\[4pt]
        \frac15
    \end{bmatrix}
    }
\]

\[
    \boxed{
    P(X_2\mid X_1=0)=
    \begin{bmatrix}
        \frac13\\[4pt]
        \frac23
    \end{bmatrix},
    \qquad
    P(X_2\mid X_1=1)=
    \begin{bmatrix}
        \frac12\\[4pt]
        \frac12
    \end{bmatrix}
    }
\]

\section*{Part (b): Gibbs Sampling for 3 Iterations}

\subsection*{Step 5: Initialization from Independent Uniform Distributions}

The question says to start with uniform distributions for \(X_1\) and \(X_2\)
independently. For a binary variable, that means
\[
    P(X_j=0)=P(X_j=1)=\frac12.
\]

Using the first two random numbers:
\begin{itemize}
    \item \(r_1=0.94\). Since \(0.94>0.5\), take
    \[
        X_1^{(0)}=1.
    \]
    \item \(r_2=0.27\). Since \(0.27\le 0.5\), take
    \[
        X_2^{(0)}=0.
    \]
\end{itemize}

So the initial state is
\[
    (X_1^{(0)},X_2^{(0)})=(1,0).
\]

\subsection*{Step 6: Iteration 1}

Sample \(X_1^{(1)}\) from \(P(X_1\mid X_2=0)\):
\[
    P(X_1=0\mid X_2=0)=\frac23.
\]
Use \(r_3=0.30\). Since \(0.30\le \frac23\),
\[
    X_1^{(1)}=0.
\]

Now sample \(X_2^{(1)}\) from \(P(X_2\mid X_1=0)\):
\[
    P(X_2=0\mid X_1=0)=\frac13.
\]
Use \(r_4=0.18\). Since \(0.18\le \frac13\),
\[
    X_2^{(1)}=0.
\]

So
\[
    (X_1^{(1)},X_2^{(1)})=(0,0).
\]

\subsection*{Step 7: Iteration 2}

Sample \(X_1^{(2)}\) from \(P(X_1\mid X_2=0)\):
\[
    P(X_1=0\mid X_2=0)=\frac23.
\]
Use \(r_5=0.75\). Since \(0.75>\frac23\),
\[
    X_1^{(2)}=1.
\]

Now sample \(X_2^{(2)}\) from \(P(X_2\mid X_1=1)\):
\[
    P(X_2=0\mid X_1=1)=\frac12.
\]
Use \(r_6=0.35\). Since \(0.35\le \frac12\),
\[
    X_2^{(2)}=0.
\]

So
\[
    (X_1^{(2)},X_2^{(2)})=(1,0).
\]

\subsection*{Step 8: Iteration 3}

Sample \(X_1^{(3)}\) from \(P(X_1\mid X_2=0)\):
\[
    P(X_1=0\mid X_2=0)=\frac23.
\]
Use \(r_7=0.61\). Since \(0.61\le \frac23\),
\[
    X_1^{(3)}=0.
\]

Now sample \(X_2^{(3)}\) from \(P(X_2\mid X_1=0)\):
\[
    P(X_2=0\mid X_1=0)=\frac13.
\]
Use \(r_8=0.42\). Since \(0.42>\frac13\),
\[
    X_2^{(3)}=1.
\]

So
\[
    (X_1^{(3)},X_2^{(3)})=(0,1).
\]

\subsection*{Answer for Part (b)}

The Gibbs samples are
\[
    \boxed{
    (X_1^{(0)},X_2^{(0)})=(1,0),\quad
    (X_1^{(1)},X_2^{(1)})=(0,0),\quad
    (X_1^{(2)},X_2^{(2)})=(1,0),\quad
    (X_1^{(3)},X_2^{(3)})=(0,1).
    }
\]

\section*{Why the Procedure Works}

Gibbs sampling does not sample from the full joint distribution directly. Instead,
it alternates:
\begin{itemize}
    \item sample \(X_1\) from \(P(X_1\mid X_2)\),
    \item then sample \(X_2\) from \(P(X_2\mid X_1)\).
\end{itemize}

These conditional updates define a Markov chain whose limiting distribution is
the target joint distribution.

\section*{Exam Shortcut}

For a small Gibbs-sampling question:
\begin{enumerate}[label=\arabic*.]
    \item Write the joint table clearly.
    \item Compute marginals.
    \item Compute the conditional tables once.
    \item Use the random numbers in order, comparing each one with the relevant
    cumulative probability.
\end{enumerate}

For binary variables, each update is just a threshold test.

\section*{Sources Used}

\begin{itemize}
    \item Assignment 4 PDF, Question 8.
    \item Local submitted Assignment 4 PDF checked for consistency of the conditionals and sample path.
    \item Chapter 7 slides on Gibbs sampling and conditional sampling.
    \item Goodfellow, Bengio, and Courville, \emph{Deep Learning}, Chapter 17 on MCMC and repeated conditional updates.
\end{itemize}

\end{document}
